\documentclass[../../main.tex]{subfiles}

\begin{document}

\section{度量空间的例子}\label{section:ch09-度量空间的例子}

\begin{definition}[度量]\label{definition:ch09-度量}
令 $X$ 为非空集. 函数 $\rho:X\times X\to\mathbf{R}$ 称为一个度量, 若对于所有属于 $X$ 的 $x$, $y$ 和 $z$,
\begin{align*}
\text{(i) }\rho(x,y)\geqslant 0;
\end{align*}
\begin{align*}
\text{(ii) }\rho(x,y)=0\text{ 当且仅当 }x=y;
\end{align*}
\begin{align*}
\text{(iii) }\rho(x,y)=\rho(y,x);
\end{align*}
\begin{align*}
\text{(iv) }\rho(x,y)\leqslant\rho(x,z)+\rho(z,y).
\end{align*}
具有一个度量的非空集合称为度量空间.
\end{definition}

我们常将度量空间记为 $(X,\rho)$. 性质 (iv) 即熟知的关于度量的三角不等式. 度量空间的一个典型例子是具有度量 $\rho(x,y)=|x-y|$ 的全体实数集 $\mathbf{R}$.

\textbf{赋范线性空间} 在 7.1 节我们把绝对值的概念推广到一般的线性空间. 回忆线性空间 $X$ 上的非负实值函数 $\|\cdot\|$ 称为范数, 若对每个 $u$, $v\in X$ 和实数 $\alpha$, 满足
\begin{align*}
\text{(i) }\|u\|=0\text{ 当且仅当 }u=0.
\end{align*}
\begin{align*}
\text{(ii) }\|u+v\|\leqslant\|u\|+\|v\|.
\end{align*}
\begin{align*}
\text{(iii) }\|\alpha u\|=|\alpha|\|u\|.
\end{align*}

我们称具有范数的线性空间为赋范线性空间. 线性空间 $X$ 上的范数 $\|\cdot\|$ 通过对所有 $x$, $y\in X$ 定义
\begin{align}
\rho(x,y)=\|x-y\|\label{align:ch09-范数诱导度量}
\end{align}
诱导了 $X$ 上的度量 $\rho$.

范数的性质 (ii) 称为关于范数的三角不等式. 它等价于其诱导的度量的三角不等式. 事实上, 对 $x$, $y$, $z\in X$, 设 $u=x-z$ 与 $v=z-y$ 且观察到
\begin{align*}
\|u+v\|\leqslant\|u\|+\|v\|\text{ 当且仅当 }\rho(x,y)\leqslant\rho(x,z)+\rho(z,y).
\end{align*}
赋范线性空间的三个重要例子是欧氏空间 $\mathbf{R}^{n}$, $L^{p}(E)$ 空间和 $C[a,b]$. 对于自然数 $n$, 考虑线性空间 $\mathbf{R}^{n}$, 其点为 $n$ 元实数组. 对于 $\mathbf{R}^{n}$ 中的 $x=(x_{1},\cdots,x_{n})$, $x$ 的欧氏范数 $\|x\|$ 定义为
\begin{align*}
\|x\|=[x_{1}^{2}+\cdots+x_{n}^{2}]^{1/2}.
\end{align*}
我们把第 7 章和第 8 章用于赋范线性空间 $L^{p}(E)$ 的研究, 其中 $1\leqslant p\leqslant\infty$ 而 $E$ 是一个实数的 Lebesgue 可测集. 对于 $1\leqslant p<\infty$, $L^{p}(E)$ 范数的三角不等式称为 Minkowski 不等式. 对于实数的闭有界区间 $[a,b]$, 考虑 $[a,b]$ 上的连续实值函数的线性空间 $C[a,b]$. $f\in C[a,b]$ 的最大值范数 $\|\cdot\|_{\max}$ 定义为
\begin{align*}
\|f\|_{\max}=\max\{|f(x)|\mid x\in[a,b]\}.
\end{align*}
最大值范数的三角不等式从实数的绝对值的三角不等式得到.

\textbf{离散度量} 对于任何非空集合 $X$, 定义离散度量 $\rho$ 为: 若 $x=y$, 则 $\rho(x,y)=0$; 而若 $x\neq y$, 则 $\rho(x,y)=1$.

\textbf{度量子空间} 对于度量空间 $(X,\rho)$, 令 $Y$ 为 $X$ 的非空子集. 则 $\rho$ 在 $Y\times Y$ 上的限制定义了 $Y$ 上的一个度量. 我们称这样的度量空间为度量子空间. 因此欧氏空间、 $L^{p}(E)$ 空间 $(1\leqslant p\leqslant\infty)$ 以及 $C[a,b]$ 的每个非空子集都是度量空间.

\textbf{度量的乘积} 对于度量空间 $(X_{1},\rho_{1})$ 和 $(X_{2},\rho_{2})$, 我们定义笛卡儿积 $X_{1}\times X_{2}$ 上的乘积度量 $\tau$ 为对于 $X_{1}\times X_{2}$ 中的 $(x_{1},x_{2})$ 和 $(y_{1},y_{2})$,
\begin{align*}
\tau((x_{1},x_{2}),(y_{1},y_{2}))=\{[\rho_{1}(x_{1},y_{1})]^{2}+[\rho_{2}(x_{2},y_{2})]^{2}\}^{1/2}.
\end{align*}
容易验证 $\tau$ 具有度量所要求的所有性质. 这种构造方式可推广到可数乘积 (见习题 10).

度量空间的一个特别有趣和有用但不是作为赋范线性空间的度量子空间直接给出的例子是, $\mathbf{R}$ 的 Lebesgue 可测子集的可测子集族上的 Nikodym 度量 (见习题 5).

在由多于一点组成的非空集合 $X$ 上存在不同的度量. 例如, 若 $X$ 是闭有界区间 $[a,b]$ 上的连续函数的非空族, 则 $X$ 是一个关于离散度量、关于最大值范数诱导的度量和关于 $L^{p}[a,b]$ $(1\leqslant p<\infty)$ 范数诱导的度量的度量空间. 以下关于集合上的度量的等价关系是有用的.

\begin{definition}[等价度量]\label{definition:ch09-等价度量}
集合 $X$ 上的两个度量 $\rho$ 和 $\sigma$ 称为等价的, 若存在正数 $c_{1}$ 和 $c_{2}$, 使得对所有 $x_{1}$, $x_{2}\in X$,
\begin{align*}
c_{1}\cdot\sigma(x_{1},x_{2})\leqslant\rho(x_{1},x_{2})\leqslant c_{2}\cdot\sigma(x_{1},x_{2}).
\end{align*}
\end{definition}

\begin{definition}[等距映射]\label{definition:ch09-等距映射}
从度量空间 $(X,\rho)$ 到度量空间 $(Y,\sigma)$ 的映射 $f$ 称为等距映射, 若它将 $X$ 映上 $Y$ 且对所有 $x_{1}$, $x_{2}\in X$,
\begin{align*}
\sigma(f(x_{1}),f(x_{2}))=\rho(x_{1},x_{2}).
\end{align*}
\end{definition}

两个度量空间称为等距的, 若存在将其中的一个映上另一个的等距映射. 等距是度量空间之间的等价关系. 从度量空间的观点看, 两个等距的度量空间是相同的. 等距映射仅相当于点的重新标记.

在集合 $X$ 上的度量 $\rho$ 的定义中, 放松条件 $\rho(x,y)=0$ 仅当 $x=y$ 有时是方便的. 当我们允许对某个 $x\neq y$ 有 $\rho(x,y)=0$ 时, 称 $\rho$ 是一个伪度量, 而 $(X,\rho)$ 是一个伪度量空间. 在这样的空间上, 定义关系 $x\cong y$, 若 $\rho(x,y)=0$. 这是一个等价关系, 它将 $X$ 分成不交的等价类族 $X/\cong$. 对等价类 $[x]$ 和 $[y]$, 定义 $\widetilde{\rho}([x],[y])=\rho(x,y)$. 容易看出这个性质定义了 $X/\cong$ 上的一个度量 $\widetilde{\rho}$. 当我们在范数的定义中对 $u\neq 0$ 允许 $\|u\|=0$ 时, 类似的考虑是适用的. 习题 5、 7、 9 以及 50 考虑了伪度量和伪范数的一些例子.

\begin{exercise}\label{exercise:ch09-01}
证明同一个集合 $X$ 上的两个度量 $\rho$ 和 $\tau$ 等价, 当且仅当存在 $c>0$ 使得对所有 $u$, $v\in X$,
\begin{align*}
\frac{1}{c}\tau(u,v)\leqslant\rho(u,v)\leqslant c\tau(u,v).
\end{align*}
\end{exercise}

\begin{exercise}\label{exercise:ch09-02}
证明以下定义了 $\mathbf{R}^{n}$ 上的等价度量:
\begin{align*}
\rho^{+}(x,y)=|x_{1}-y_{1}|+\cdots+|x_{n}-y_{n}|,
\end{align*}
\begin{align*}
\rho^{-}(x,y)=\max\{|x_{1}-y_{1}|,\cdots,|x_{n}-y_{n}|\}.
\end{align*}
\end{exercise}

\begin{exercise}\label{exercise:ch09-03}
找到 $\mathbf{R}^{n}$ 上与前一个习题中的度量都不等价的度量.
\end{exercise}

\begin{exercise}\label{exercise:ch09-04}
对一个闭有界区间 $[a,b]$, 考虑 $[a,b]$ 上的连续实值函数集合 $X=C[a,b]$. 证明由最大值范数诱导的度量与由 $L^{1}[a,b]$ 范数诱导的度量不等价.
\end{exercise}

\begin{exercise}\label{exercise:ch09-05}
(Nikodym 度量) 令 $E$ 为具有有限测度的实数的 Lebesgue 可测集, $X$ 是 $E$ 的 Lebesgue 可测子集的集合, 而 $m$ 是 Lebesgue 测度. 对 $A$, $B\in X$, 定义 $\rho(A,B)=m(A\Delta B)$, 其中 $A\Delta B=[A\sim B]\bigcup[B\sim A]$ 是集合 $A$ 和 $B$ 的对称差. 证明这是 $X$ 上的伪度量. 定义两个可测集为等价, 若它们的对称差测度为零. 证明 $\rho$ 诱导了等价类族上的度量. 最后, 证明对于 $A$, $B\in X$,
\begin{align*}
\rho(A,B)=\int_{E}|\chi_{A}-\chi_{B}|,
\end{align*}
其中 $\chi_{A}$ 和 $\chi_{B}$ 分别是 $A$ 和 $B$ 的特征函数.
\end{exercise}

\begin{exercise}\label{exercise:ch09-06}
证明对于 $a$, $b$, $c\geqslant 0$, 若 $a\leqslant b+c$, 则
\begin{align*}
\frac{a}{1+a}\leqslant\frac{b}{1+b}+\frac{c}{1+c}.
\end{align*}
\end{exercise}

\begin{exercise}\label{exercise:ch09-07}
令 $E$ 为实数的具有有限测度的 Lebesgue 可测集, 而 $X$ 为 $E$ 上的 Lebesgue 可测实值函数的集合. 对于 $f$, $g\in X$, 定义
\begin{align*}
\rho(f,g)=\int_{E}\frac{|f-g|}{1+|f-g|}.
\end{align*}
用前一个习题证明这是 $X$ 上的一个伪度量. 定义两个可测函数为等价, 若它们在 $E$ 上几乎处处相等. 证明 $\rho$ 诱导了等价类族上的度量.
\end{exercise}

\begin{exercise}\label{exercise:ch09-08}
对于 $0<p<1$, 证明: 对所有 $a$, $b\geqslant 0$, $(a+b)^{p}\leqslant a^{p}+b^{p}$.
\end{exercise}

\begin{exercise}\label{exercise:ch09-09}
对实数的 Lebesgue 可测集 $E$, $0<p<1$, $g$ 和 $h$ 是 $E$ 上的 $p$ 次幂可积的 Lebesgue 可测函数, 定义
\begin{align*}
\rho_{p}(h,g)=\int_{E}|g-h|^{p}.
\end{align*}
利用前一个习题证明这是 $E$ 上的 $p$ 次幂可积的 Lebesgue 可测函数族上的一个伪度量. 若两个这样的函数在 $E$ 上几乎处处相等, 定义它们为等价的. 证明 $\rho_{p}(\cdot,\cdot)$ 诱导了等价类族上的度量.
\end{exercise}

\begin{exercise}\label{exercise:ch09-10}
令 $\{(X_{n},\rho_{n})\}_{n=1}^{\infty}$ 为度量空间的可数族. 利用习题 6 证明 $\rho_{*}$ 定义了笛卡儿积 $\prod_{n=1}^{\infty}X_{n}$ 上的度量, 其中对于 $\prod_{n=1}^{\infty}X_{n}$ 中的点 $x=\{x_{n}\}$ 和 $y=\{y_{n}\}$,
\begin{align*}
\rho_{*}(x,y)=\sum_{n=1}^{\infty}\frac{1}{2^{n}}\cdot\frac{\rho_{n}(x_{n},y_{n})}{1+\rho_{n}(x_{n},y_{n})}.
\end{align*}
\end{exercise}

\begin{exercise}\label{exercise:ch09-11}
令 $(X,\rho)$ 为度量空间, 而 $A$ 是任何集, 使得存在将 $A$ 映上 $X$ 的一一映射 $f$. 证明 $A$ 上存在唯一的度量使得 $f$ 是度量空间之间的等距映射.
\end{exercise}

\begin{exercise}\label{exercise:ch09-12}
证明关于欧氏空间 $\mathbf{R}^{n}$ 的三角不等式从关于 $L^{2}[0,1]$ 的三角不等式得到.
\end{exercise}

\end{document}